%!TEX root = ../report.tex
\documentclass[../report.tex]{subfiles}

\begin{document}
    \section{Introduction}
    \label{sec:introduction}

    \subsection{Motivation}
    \label{sec:introduction:motivation}

    Robots have the potential to become important partners in human environments, assisting with tasks at home or in complex collaborative work \cite{Chetouani2021, Akalin2021}. However, enabling robots to operate well in open-world environments---with unknown tasks, new objects, and dynamic human interactions---remains a central challenge. Robots often lack the common-sense contextual understanding that humans use effortlessly to generalise skills and interpret underspecified instructions.

    A promising direction to address this challenge is Interactive Reinforcement Learning (IRL), which incorporates human feedback directly into the robot's learning loop \cite{Lin2020, BradleyKnox2008}. By learning from evaluative or instructive feedback, IRL agents can adapt to complex, real-world constraints more effectively than autonomous RL agents that rely solely on engineered reward functions, while also improving sample efficiency and aligning behaviour with human preferences \cite{Lin2020}.

    Despite its potential, a significant bottleneck in IRL is the interface through which feedback is delivered. Most existing systems rely on low-bandwidth, explicit channels such as keyboard presses or button clicks \cite{Faulkner2020, Loftin2014}. These methods are unnatural for the human teacher and fail to capture the rich details of human communication---gesture, gaze, tone of voice---creating an \emph{expressiveness gap} that limits the quality of information transferred from human to robot \cite{9794595, Lin2020a}.

    At the same time, advances in consumer-grade Virtual Reality (VR) technology offer new opportunities for human-computer interaction. Modern VR headsets provide dense multimodal sensing: head pose, hand gestures, eye gaze, and voice input within a single, immersive platform \cite{Winkler2022, Luo2024}. Despite this potential, VR and IRL have mostly developed in parallel. VR has been used for collecting demonstration data \cite{10904309, Mosbach2022} and, more recently, for preference-based learning \cite{Heuvel2025}, but its use for real-time scalar feedback within an IRL framework remains largely unexplored.

    \subsection{Problem Statement}
    \label{sec:introduction:problem_statement}

    The combination of the expressiveness gap in existing IRL interfaces, the limited exploitation of VR's real-time sensing capabilities for feedback, and the unresolved challenge of multimodal signal fusion defines the core problem addressed in this work. Specifically:

    \begin{itemize}
        \item Current IRL feedback interfaces are narrow-bandwidth (typically binary key presses) and unnatural, limiting the quality of human supervision.
        \item VR headsets provide rich, synchronised multimodal signals (head pose, gaze, hand gestures, voice) that have not been systematically leveraged for real-time IRL.
        \item It is unclear which combinations of VR signals are most effective as reward proxies, or how to fuse them robustly in the presence of noise and conflicts.
        \item Standard on-policy RL algorithms (e.g., PPO) require large batches of environment interactions per gradient update, making them impractical when the human feedback budget is limited to a small number of interaction steps.
    \end{itemize}

    This work addresses these gaps through the following research questions:
    \begin{enumerate}
        \item \textbf{Feasibility:} Can multimodal feedback from a VR headset be effectively translated into a robust and informative reward signal for an RL agent in a greeting interaction task?
        \item \textbf{Performance Comparison:} How does the learning efficiency and final performance of an agent fine-tuned with VR feedback compare to one trained with traditional keyboard input or an autonomous reward function?
        \item \textbf{User Experience:} What are the qualitative differences in the teaching experience between a VR interface and a keyboard interface for providing IRL feedback?
    \end{enumerate}

    \subsection{Proposed Approach}
    \label{sec:introduction:proposed_approach}

    We develop a proof-of-concept IRL system built around a social robot greeting task implemented in the Unity game engine. A Pepper robot agent must learn to select the contextually appropriate greeting behaviour---DoNothing, Talk, Look, Wave, or HandShake---in response to cues provided by an approaching human. Three conditions are designed:

    \begin{enumerate}
        \item \textbf{Autonomous RL Baseline:} The agent learns from scratch using PPO with a hand-crafted reward function (+1.0 for correct actions, $-$0.2 for incorrect actions, $-$0.0001 per-step penalty) for $\sim$150k steps. This establishes performance benchmarks and exposes the Talk action suppression that motivates the fine-tuning conditions.

        \item \textbf{Keyboard IRL:} The best PPO checkpoint is loaded into the COACH agent. A human teacher provides scalar feedback (+1/$-$1) via keyboard arrow keys. COACH performs immediate per-event policy gradient updates suited to the limited human interaction budget; the session ran for 164 steps.

        \item \textbf{VR IRL:} The same PPO checkpoint is independently loaded into a second COACH agent. A human teacher wearing a VR headset provides feedback through multiple modalities: controller buttons (explicit $\pm$1), head gestures (nod for positive, shake for negative), and voice commands (``good''/``bad''). This condition was fully implemented but could not be executed due to a physical body incompatibility between the NPC rig and the VR person controller.
    \end{enumerate}

    Conditions (2) and (3) both branch from the same PPO baseline checkpoint independently---neither depends on the other. This allows a direct comparison of feedback modality (keyboard vs VR) against each other and against the autonomous baseline. Body-signal observations are normalised relative to each person's standing height, ensuring scale-invariant features across the NPC training rig and VR participants of differing heights.

    \subsection{Contributions}
    \label{sec:introduction:contributions}

    This report makes the following specific contributions:

    \begin{enumerate}
        \item \textbf{Complete System Architecture:} Integration of Unity ML-Agents with PyTorch PPO and COACH, including a 12-dimensional observation space with body-signal features (wrist height, wrist-to-core distance, body orientation, gaze direction) that replace explicit task IDs.

        \item \textbf{Hyperparameter Characterisation:} Systematic sweep of 10 configurations varying learning rate ($10^{-4}$ to $10^{-3}$), network capacity (64, 128, 256 units), entropy schedule, rollout length, and PPO clip parameter, providing guidance for future social greeting tasks.

        \item \textbf{Performance Benchmarking:} Quantitative results establishing final training accuracy (66.4\%), test accuracy (74.14\%), sample efficiency ($\sim$9k steps to 80\% accuracy), and policy stability ($\sigma = 0.72$) for the autonomous condition as reference values for IRL comparison.

        \item \textbf{COACH Fine-Tuning Pipeline:} Implementation of COACH as a targeted fine-tuning mechanism that corrects the Talk action suppression (0\% distribution) introduced by body-signal feature overlap with the Look condition, using only 164 human feedback events loaded from the pre-trained PPO checkpoint.

        \item \textbf{Body-Relative Observation Normalisation:} A scale-invariant normalisation scheme in which wrist height and wrist-to-core distance are expressed as fractions of the person's own standing height, enabling consistent feature representations across NPC training and live VR human participants.

        \item \textbf{Multimodal VR Feedback Interface:} A priority-based fusion of controller buttons, head gestures, and voice commands into a scalar reward signal, with per-modality reward magnitudes and timeout gating to prevent signal duplication.

        \item \textbf{Reproducible Configuration:} Complete YAML configuration files (config1.yaml--config10.yaml and coach\_config.yaml) and training code for replication and extension to further IRL conditions.
    \end{enumerate}

    \subsection{Report Structure}
    \label{sec:introduction:structure}

    Section~\ref{sec:related_work} reviews related work in IRL, VR-based robot learning, and avatar motion estimation. Section~\ref{sec:background} introduces background concepts in reinforcement learning, PPO, IRL, and the COACH algorithm. Section~\ref{sec:methodology} describes the complete system design, including the training conditions, observation/action spaces, reward function, body-relative normalisation, and COACH fine-tuning methodology. Section~\ref{sec:evaluation} presents experimental results from the hyperparameter sweep, definitive PPO run, and keyboard COACH fine-tuning evaluation. Section~\ref{sec:conclusions} summarises findings and discusses implications for VR-based IRL research.

\end{document}